% Auto-generated by dump_nf_bcond_table.py
% Requires: booktabs, longtable, amsmath, amssymb

\renewcommand{\arraystretch}{1.4}

\section*{Asymptotic freedom beyond the UV}

For each non-Veneziano class we study the effect of adding $N_f$ fundamental chiral multiplets to one or both gauge nodes of the $\mathrm{SU}(N)\times\mathrm{SU}(N)$ quiver. The asymptotic-freedom condition $b_0^{(a)} > 0$ on each node must hold not only at the free UV fixed point but also at points where one coupling has flowed to its interacting fixed point while the other remains free. At such points the bifundamental R-charges are no longer free-field values but are fixed by the interacting node's dynamics, modifying the one-loop coefficient and tightening the upper bound on $N_f$. We label three relevant points in the two-dimensional $(g_1, g_2)$ coupling space as follows:
\begin{itemize}
\item $\mathcal{O} = (g_1=0,\,g_2=0)$: the free UV fixed point. The one-loop AF condition $b_0^{(a)} > 0$ at $\mathcal{O}$ gives the standard asymptotic-freedom upper bound
  $N_f < N_f^{\mathcal{O},(a)} \equiv \alpha_{\mathcal{O}}^{(a)}\,N + \beta^{(a)},$
  where $\alpha_{\mathcal{O}}^{(a)}$ is determined by the rank-2 and bifundamental content of node $a$.
\item $\mathcal{B} = (g_1=0,\,g_2=g_2^*)$: node 2 is at its interacting fixed point, node 1 is free. The R-charges of the bifundamentals are now fixed by node 2's fixed point. The condition $b_0^{(1)}|_{\mathcal{B}} > 0$ gives
  $N_f < N_f^{\mathcal{B},(1)} \equiv \alpha_{\mathcal{B}}^{(1)}\,N + \beta^{(1)}.$
\item $\mathcal{A} = (g_1=g_1^*,\,g_2=0)$: node 1 is at its interacting fixed point, node 2 is free. The condition $b_0^{(2)}|_{\mathcal{A}} > 0$ gives
  $N_f < N_f^{\mathcal{A},(2)} \equiv \alpha_{\mathcal{A}}^{(2)}\,N + \beta^{(2)}.$
\end{itemize}

The theory flows to a non-trivial IR fixed point only if $N_f$ satisfies all three bounds simultaneously. The coefficient $\alpha$ of $N$ is fixed by the morphology $(N_{r2}^{(1)}, N_{r2}^{(2)}, N_{\rm bif})$; the constant $\beta$ depends on the S/A character of the rank-2 tensors but does not affect the large-$N$ scaling and is not tabulated here. Each table below lists the coefficients $\alpha$ for all morphologies within a class; negative values of $\alpha_{\mathcal{B}}$ or $\alpha_{\mathcal{A}}$ indicate that the corresponding boundary condition is automatically satisfied for large $N$ and imposes no constraint on the conformal window.

\subsection*{Class~1: $a/N^2 = 0$\quad (6 theories)}

This class has $a/N^2 = 0$, the lowest possible value, corresponding to a fixed point at which all matter R-charges vanish identically. The $\mathcal{B}$-condition and the $\mathcal{O}$-condition therefore differ only in the rank-2 tensor contribution, yielding a simple uniform compression of the conformal window by a factor of two.

\textit{Morphology $(\tfrac{1}{2},\,\tfrac{1}{2},\,1)$.}\;With one rank-2 tensor on each node and a single bifundamental, the $\mathcal{O}$-point beta function gives $\alpha_{\mathcal{O}} = 2$ on both nodes. At $a/N^2 = 0$ all matter R-charges vanish at the fixed point, so the bifundamental enters the $\mathcal{B}$-point (and $\mathcal{A}$-point) beta function with R-charge zero, exactly as at $\mathcal{O}$. The difference between the two bounds arises solely from the rank-2 tensor contribution: at $\mathcal{B}$ the rank-2 tensor on node 1 is free ($R = 0$) while the bifundamental contribution from node 2 is unchanged, halving the net coefficient. The result is a uniform compression $\alpha_{\mathcal{O}} = 2 \to \alpha_{\mathcal{B}} = 1$ on both nodes.

\begin{longtable}{l|cc|cc}
\caption{Coefficients $\alpha$ of $N$ in the upper bounds $N_f < \alpha\,N + \beta$ at each boundary point, for all morphologies of class~1. The constant $\beta$ (not shown) depends on the S/A character of the rank-2 matter. Negative $\alpha_{\mathcal{B}}$ or $\alpha_{\mathcal{A}}$ means the boundary condition is automatically satisfied at large $N$; `conf.' means the node is already at the conformal threshold at $\mathcal{O}$.}\label{tab:nf_cl1}\\
\toprule
$(N_{r2}^{(1)},\,N_{r2}^{(2)},\,N_{\rm bif})$ & $\alpha_{\mathcal{O}}^{(1)}$ & $\alpha_{\mathcal{B}}^{(1)}$ & $\alpha_{\mathcal{O}}^{(2)}$ & $\alpha_{\mathcal{A}}^{(2)}$ \\
\midrule
\endhead
$(\tfrac{1}{2},\,\tfrac{1}{2},\,1)$ & $2$ & $1$ & $2$ & $1$ \\
\bottomrule
\end{longtable}

\subsection*{Class~2: $a/N^2 = -\frac{297}{32} + 135\sqrt{5}/32$\quad (32 theories)}

The fixed point involves irrational R-charges proportional to $\sqrt{5}$. The class has two morphologies related by node exchange. In both, the same shift $\Delta\alpha \approx 0.781$ tightens both boundary conditions relative to $\mathcal{O}$. The binding constraint comes from the heavier node (larger $N_{r2}$, smaller $\alpha_{\mathcal{O}}$), whose tightened coefficient $\approx 0.219$ controls the conformal window.

\textit{Morphology $(\tfrac{1}{2},\,\tfrac{3}{2},\,1)$.}\;Node 1 carries one rank-2 tensor ($\alpha_{\mathcal{O}}^{(1)} = 2$) and node 2 carries three ($\alpha_{\mathcal{O}}^{(2)} = 1$). At the interacting fixed point the bifundamental acquires R-charge $R_{\rm bif} \approx 0.146 < \tfrac{2}{3}$, giving a negative anomalous dimension $\gamma_{\rm bif} = 3R_{\rm bif} - 2 \approx -1.562$. This increased depletion of the one-loop coefficient tightens both boundary conditions by the same amount $\Delta\alpha = \tfrac{1}{2}|\gamma_{\rm bif}| \approx 0.781$: $\alpha_{\mathcal{B}}^{(1)} \approx 1.219 < \alpha_{\mathcal{O}}^{(1)} = 2$ and $\alpha_{\mathcal{A}}^{(2)} \approx 0.219 < \alpha_{\mathcal{O}}^{(2)} = 1$. The binding constraint on the conformal window is $\alpha_{\mathcal{A}}^{(2)} \approx 0.219$.

\textit{Morphology $(\tfrac{3}{2},\,\tfrac{1}{2},\,1)$.}\;The node-exchange mirror of the previous morphology. The same $\Delta\alpha \approx 0.781$ tightening applies: $\alpha_{\mathcal{B}}^{(1)} \approx 0.219 < \alpha_{\mathcal{O}}^{(1)} = 1$ and $\alpha_{\mathcal{A}}^{(2)} \approx 1.219 < \alpha_{\mathcal{O}}^{(2)} = 2$. The binding constraint is now $\alpha_{\mathcal{B}}^{(1)} \approx 0.219$.

\begin{longtable}{l|cc|cc}
\caption{Coefficients $\alpha$ of $N$ in the upper bounds $N_f < \alpha\,N + \beta$ at each boundary point, for all morphologies of class~2. The constant $\beta$ (not shown) depends on the S/A character of the rank-2 matter. Negative $\alpha_{\mathcal{B}}$ or $\alpha_{\mathcal{A}}$ means the boundary condition is automatically satisfied at large $N$; `conf.' means the node is already at the conformal threshold at $\mathcal{O}$.}\label{tab:nf_cl2}\\
\toprule
$(N_{r2}^{(1)},\,N_{r2}^{(2)},\,N_{\rm bif})$ & $\alpha_{\mathcal{O}}^{(1)}$ & $\alpha_{\mathcal{B}}^{(1)}$ & $\alpha_{\mathcal{O}}^{(2)}$ & $\alpha_{\mathcal{A}}^{(2)}$ \\
\midrule
\endhead
$(\tfrac{1}{2},\,\tfrac{3}{2},\,1)$ & $2$ & $1.2188$ & $1$ & $0.2188$ \\
$(\tfrac{3}{2},\,\tfrac{1}{2},\,1)$ & $1$ & $0.2188$ & $2$ & $1.2188$ \\
\bottomrule
\end{longtable}

\subsection*{Class~4: $a/N^2 = -\frac{229}{32} + 65\sqrt{13}/32$\quad (12 theories)}

This class has irrational R-charges proportional to $\sqrt{13}$ and a single asymmetric morphology. The heavier node saturates the $\mathcal{O}$-point beta function ($\alpha_{\mathcal{O}}^{(1)} = 0$): the same tightening shift $\Delta\alpha \approx 0.757$ pushes $\alpha_{\mathcal{B}}^{(1)} \approx -0.757 < 0$, automatically satisfying $b_0^{(1)}|_{\mathcal{B}} < 0$ for large $N$. The conformal window is controlled by the lighter node's tightened coefficient $\alpha_{\mathcal{A}}^{(2)} \approx 1.243$.

\textit{Morphology $(\tfrac{5}{2},\,\tfrac{1}{2},\,1)$.}\;Node 1 carries five half-units of rank-2 matter ($N_{r2}^{(1)} = \tfrac{5}{2}$), saturating the $\mathcal{O}$-point beta function: $\alpha_{\mathcal{O}}^{(1)} = 0$. The bifundamental R-charge at the $\sqrt{13}$ fixed point is $R_{\rm bif} \approx 0.162$, giving $\Delta\alpha \approx 0.757$. Since $\alpha_{\mathcal{O}}^{(1)} = 0$, the tightening pushes node 1's coefficient below zero: $\alpha_{\mathcal{B}}^{(1)} \approx -0.757 < 0$, so $b_0^{(1)}|_{\mathcal{B}} < 0$ is automatically satisfied for any $N_f \geq 0$ at large $N$. Node 2, with one rank-2 tensor ($\alpha_{\mathcal{O}}^{(2)} = 2$), is tightened to $\alpha_{\mathcal{A}}^{(2)} \approx 1.243$: the conformal window is controlled by $b_0^{(2)}|_{\mathcal{A}}$.

\begin{longtable}{l|cc|cc}
\caption{Coefficients $\alpha$ of $N$ in the upper bounds $N_f < \alpha\,N + \beta$ at each boundary point, for all morphologies of class~4. The constant $\beta$ (not shown) depends on the S/A character of the rank-2 matter. Negative $\alpha_{\mathcal{B}}$ or $\alpha_{\mathcal{A}}$ means the boundary condition is automatically satisfied at large $N$; `conf.' means the node is already at the conformal threshold at $\mathcal{O}$.}\label{tab:nf_cl4}\\
\toprule
$(N_{r2}^{(1)},\,N_{r2}^{(2)},\,N_{\rm bif})$ & $\alpha_{\mathcal{O}}^{(1)}$ & $\alpha_{\mathcal{B}}^{(1)}$ & $\alpha_{\mathcal{O}}^{(2)}$ & $\alpha_{\mathcal{A}}^{(2)}$ \\
\midrule
\endhead
$(\tfrac{5}{2},\,\tfrac{1}{2},\,1)$ & $\mathrm{conf.}$ & $-0.7569$ & $2$ & $1.2431$ \\
\bottomrule
\end{longtable}

\subsection*{Class~9: $a/N^2 = \frac{27}{64}$\quad (139 theories)}

This is the class with the simplest rational structure. All morphologies share $a/N^2 = 27/64$ and have $\alpha_{\mathcal{O}} = 1$ on both nodes. The $\mathcal{B}$- (and $\mathcal{A}$-) point coefficient follows the exact rational formula\vspace{-2pt} 
$$\alpha_{\mathcal{B}} = 1 - \frac{N_{\rm bif}}{4}$$
for the symmetric morphologies ($N_{r2}^{(1)} = N_{r2}^{(2)}$). The formula counts how much the four bifundamental R-charges --- each equal to $\tfrac{1}{2}$ at the $a/N^2 = 27/64$ fixed point --- deplete the boundary beta function relative to the free-field value.

\textit{Morphology $(0,\,0,\,4)$.}\;Both nodes carry no rank-2 tensors and are connected by four bifundamentals. At $\mathcal{O}$, each node has $b_0 = N$ giving $\alpha_{\mathcal{O}} = 1$. At $\mathcal{B}$, the four bifundamentals carry their fixed-point R-charges from node 2; since $N_{\rm bif} = 4$ saturates the formula $\alpha_{\mathcal{B}} = 1 - N_{\rm bif}/4$, the coefficient collapses to zero. The conformal window for additional fundamentals is thus $N$-independent at both boundaries, anticipating the behaviour of class~44.

\textit{Morphology $(0,\,1,\,4)$.}\;Node 1 has no rank-2 tensors ($\alpha_{\mathcal{O}}^{(1)} = 1$) while node 2 carries one adjoint unit ($N_{r2}^{(2)} = 1$, $\alpha_{\mathcal{O}}^{(2)} = 0$, UV threshold). With $\Delta\alpha = N_{\rm bif}/4 = 1$: node 1's coefficient drops to $\alpha_{\mathcal{B}}^{(1)} = 0$ (conformal threshold at $\mathcal{B}$), and node 2's UV-saturated condition is pushed to $\alpha_{\mathcal{A}}^{(2)} = -1 < 0$, automatically satisfied. No fundamental deformation with $N_f \sim N$ is possible for either node.

\textit{Morphology $(\tfrac{1}{2},\,\tfrac{1}{2},\,3)$.}\;One rank-2 tensor on each node and three bifundamentals. The formula $\alpha_{\mathcal{B}} = 1 - N_{\rm bif}/4 = \tfrac{1}{4}$ applies to both nodes by symmetry. The conformal window is compressed to one-quarter of the $\mathcal{O}$-point estimate.

\textit{Morphology $(1,\,1,\,2)$.}\;Two rank-2 equivalent units per node and two bifundamentals. The formula gives $\alpha_{\mathcal{B}} = 1 - 2/4 = \tfrac{1}{2}$ on both nodes. The conformal window is compressed to one-half of the $\mathcal{O}$-point estimate.

\textit{Morphology $(\tfrac{3}{2},\,\tfrac{3}{2},\,1)$.}\;Three rank-2 equivalent units per node and one bifundamental. The formula gives $\alpha_{\mathcal{B}} = 1 - 1/4 = \tfrac{3}{4}$ on both nodes: the mildest compression within class~9, consistent with the single bifundamental carrying only a small fixed-point R-charge.

\begin{longtable}{l|cc|cc}
\caption{Coefficients $\alpha$ of $N$ in the upper bounds $N_f < \alpha\,N + \beta$ at each boundary point, for all morphologies of class~9. The constant $\beta$ (not shown) depends on the S/A character of the rank-2 matter. Negative $\alpha_{\mathcal{B}}$ or $\alpha_{\mathcal{A}}$ means the boundary condition is automatically satisfied at large $N$; `conf.' means the node is already at the conformal threshold at $\mathcal{O}$.}\label{tab:nf_cl9}\\
\toprule
$(N_{r2}^{(1)},\,N_{r2}^{(2)},\,N_{\rm bif})$ & $\alpha_{\mathcal{O}}^{(1)}$ & $\alpha_{\mathcal{B}}^{(1)}$ & $\alpha_{\mathcal{O}}^{(2)}$ & $\alpha_{\mathcal{A}}^{(2)}$ \\
\midrule
\endhead
$(0,\,0,\,4)$ & $1$ & $\mathrm{conf.}$ & $1$ & $\mathrm{conf.}$ \\
$(0,\,1,\,4)$ & $1$ & $\mathrm{conf.}$ & $\mathrm{conf.}$ & $-1$ \\
$(\tfrac{1}{2},\,\tfrac{1}{2},\,3)$ & $1$ & $\frac{1}{4}$ & $1$ & $\frac{1}{4}$ \\
$(1,\,1,\,2)$ & $1$ & $\frac{1}{2}$ & $1$ & $\frac{1}{2}$ \\
$(\tfrac{3}{2},\,\tfrac{3}{2},\,1)$ & $1$ & $\frac{3}{4}$ & $1$ & $\frac{3}{4}$ \\
\bottomrule
\end{longtable}

\subsection*{Class~16: $a/N^2 = \frac{3}{32} + 5\sqrt{5}/32$\quad (8 theories)}

Irrational R-charges proportional to $\sqrt{5}$, a single asymmetric morphology with three bifundamentals. The structure is analogous to class~4: the UV-saturated node's condition is automatically satisfied at $\mathcal{B}$, and the lighter node's condition tightens to $\alpha_{\mathcal{A}}^{(2)} \approx 0.427$ at $\mathcal{A}$.

\textit{Morphology $(\tfrac{3}{2},\,\tfrac{1}{2},\,3)$.}\;Node 1 carries $\tfrac{3}{2}$ rank-2 equivalent units with three bifundamentals ($\alpha_{\mathcal{O}}^{(1)} = 0$, at the $\mathcal{O}$-point threshold). $\Delta\alpha = \tfrac{3}{2}|\gamma_{\rm bif}| \approx 0.573$: $\alpha_{\mathcal{B}}^{(1)} \approx -0.573 < 0$ (automatically satisfying $b_0^{(1)}|_{\mathcal{B}} < 0$ for large $N$), and node 2's condition tightens to $\alpha_{\mathcal{A}}^{(2)} \approx 0.427 < \alpha_{\mathcal{O}}^{(2)} = 1$. The conformal window is entirely controlled by $b_0^{(2)}|_{\mathcal{A}}$.

\begin{longtable}{l|cc|cc}
\caption{Coefficients $\alpha$ of $N$ in the upper bounds $N_f < \alpha\,N + \beta$ at each boundary point, for all morphologies of class~16. The constant $\beta$ (not shown) depends on the S/A character of the rank-2 matter. Negative $\alpha_{\mathcal{B}}$ or $\alpha_{\mathcal{A}}$ means the boundary condition is automatically satisfied at large $N$; `conf.' means the node is already at the conformal threshold at $\mathcal{O}$.}\label{tab:nf_cl16}\\
\toprule
$(N_{r2}^{(1)},\,N_{r2}^{(2)},\,N_{\rm bif})$ & $\alpha_{\mathcal{O}}^{(1)}$ & $\alpha_{\mathcal{B}}^{(1)}$ & $\alpha_{\mathcal{O}}^{(2)}$ & $\alpha_{\mathcal{A}}^{(2)}$ \\
\midrule
\endhead
$(\tfrac{3}{2},\,\tfrac{1}{2},\,3)$ & $\mathrm{conf.}$ & $-0.5730$ & $1$ & $0.4270$ \\
\bottomrule
\end{longtable}

\subsection*{Class~21: $a/N^2 = -\frac{7}{2} + 5\sqrt{10}/4$\quad (77 theories)}

Irrational R-charges proportional to $\sqrt{10}$, two morphologies related by node exchange. In both, the UV-saturated node's boundary condition is automatically satisfied (pushed to $\alpha < 0$), and the conformal window is controlled by the other node's tightened coefficient.

\textit{Morphology $(1,\,2,\,2)$.}\;Node 1 carries one rank-2 equivalent unit ($\alpha_{\mathcal{O}}^{(1)} = 1$) and node 2 carries two ($\alpha_{\mathcal{O}}^{(2)} = 0$, at the threshold). $\Delta\alpha = N_{\rm bif}|\gamma_{\rm bif}|/2 \approx 0.325$: node 1's condition tightens to $\alpha_{\mathcal{B}}^{(1)} \approx 0.675 < \alpha_{\mathcal{O}}^{(1)} = 1$, and node 2's UV-saturated condition is pushed to $\alpha_{\mathcal{A}}^{(2)} \approx -0.325 < 0$, automatically satisfied. The conformal window is controlled by $b_0^{(1)}|_{\mathcal{B}}$.

\textit{Morphology $(2,\,1,\,2)$.}\;The mirror of the previous morphology with nodes 1 and 2 exchanged. Node 1 ($N_{r2}^{(1)} = 2$, $\alpha_{\mathcal{O}}^{(1)} = 0$) is at the $\mathcal{O}$-point threshold: $\alpha_{\mathcal{B}}^{(1)} \approx -0.325 < 0$, automatically satisfied. Node 2's condition tightens to $\alpha_{\mathcal{A}}^{(2)} \approx 0.675 < \alpha_{\mathcal{O}}^{(2)} = 1$: the conformal window is controlled by $b_0^{(2)}|_{\mathcal{A}}$.

\begin{longtable}{l|cc|cc}
\caption{Coefficients $\alpha$ of $N$ in the upper bounds $N_f < \alpha\,N + \beta$ at each boundary point, for all morphologies of class~21. The constant $\beta$ (not shown) depends on the S/A character of the rank-2 matter. Negative $\alpha_{\mathcal{B}}$ or $\alpha_{\mathcal{A}}$ means the boundary condition is automatically satisfied at large $N$; `conf.' means the node is already at the conformal threshold at $\mathcal{O}$.}\label{tab:nf_cl21}\\
\toprule
$(N_{r2}^{(1)},\,N_{r2}^{(2)},\,N_{\rm bif})$ & $\alpha_{\mathcal{O}}^{(1)}$ & $\alpha_{\mathcal{B}}^{(1)}$ & $\alpha_{\mathcal{O}}^{(2)}$ & $\alpha_{\mathcal{A}}^{(2)}$ \\
\midrule
\endhead
$(1,\,2,\,2)$ & $1$ & $0.6754$ & $\mathrm{conf.}$ & $-0.3246$ \\
$(2,\,1,\,2)$ & $\mathrm{conf.}$ & $-0.3246$ & $1$ & $0.6754$ \\
\bottomrule
\end{longtable}

\subsection*{Class~23: $a/N^2 = 2295\sqrt{17}/86528 + \frac{30159}{86528}$\quad (48 theories)}

The highest-$a$ irrational class, with R-charges proportional to $\sqrt{17}$. Two morphologies related by node exchange. The tightened boundary coefficient $\alpha \approx 0.853$ is the largest among all non-Veneziano irrational classes (at $\mathcal{B}$ or $\mathcal{A}$ depending on the morphology), approaching unity as $a/N^2 \to \tfrac{1}{2}$.

\textit{Morphology $(\tfrac{3}{2},\,\tfrac{5}{2},\,1)$.}\;Node 1 carries $\tfrac{3}{2}$ rank-2 equivalent units ($\alpha_{\mathcal{O}}^{(1)} = 1$) and node 2 carries $\tfrac{5}{2}$ ($\alpha_{\mathcal{O}}^{(2)} = 0$, at the threshold). $\Delta\alpha = \tfrac{1}{2}|\gamma_{\rm bif}| \approx 0.147$: node 1's condition tightens to $\alpha_{\mathcal{B}}^{(1)} \approx 0.853 < \alpha_{\mathcal{O}}^{(1)} = 1$, the largest tightened coefficient among all irrational classes, and node 2's UV-saturated condition is pushed to $\alpha_{\mathcal{A}}^{(2)} \approx -0.147 < 0$, automatically satisfied. The conformal window is controlled by $b_0^{(1)}|_{\mathcal{B}}$.

\textit{Morphology $(\tfrac{5}{2},\,\tfrac{3}{2},\,1)$.}\;The mirror morphology: node 1 carries $\tfrac{5}{2}$ rank-2 equivalent units ($\alpha_{\mathcal{O}}^{(1)} = 0$, at threshold): $\alpha_{\mathcal{B}}^{(1)} \approx -0.147 < 0$, automatically satisfied. Node 2 ($\tfrac{3}{2}$ rank-2 units, $\alpha_{\mathcal{O}}^{(2)} = 1$) sees its condition tighten to $\alpha_{\mathcal{A}}^{(2)} \approx 0.853$: the conformal window is controlled by $b_0^{(2)}|_{\mathcal{A}}$.

\begin{longtable}{l|cc|cc}
\caption{Coefficients $\alpha$ of $N$ in the upper bounds $N_f < \alpha\,N + \beta$ at each boundary point, for all morphologies of class~23. The constant $\beta$ (not shown) depends on the S/A character of the rank-2 matter. Negative $\alpha_{\mathcal{B}}$ or $\alpha_{\mathcal{A}}$ means the boundary condition is automatically satisfied at large $N$; `conf.' means the node is already at the conformal threshold at $\mathcal{O}$.}\label{tab:nf_cl23}\\
\toprule
$(N_{r2}^{(1)},\,N_{r2}^{(2)},\,N_{\rm bif})$ & $\alpha_{\mathcal{O}}^{(1)}$ & $\alpha_{\mathcal{B}}^{(1)}$ & $\alpha_{\mathcal{O}}^{(2)}$ & $\alpha_{\mathcal{A}}^{(2)}$ \\
\midrule
\endhead
$(\tfrac{3}{2},\,\tfrac{5}{2},\,1)$ & $1$ & $0.8532$ & $\mathrm{conf.}$ & $-0.1468$ \\
$(\tfrac{5}{2},\,\tfrac{3}{2},\,1)$ & $\mathrm{conf.}$ & $-0.1468$ & $1$ & $0.8532$ \\
\bottomrule
\end{longtable}

\subsection*{Class~44: $a/N^2 = \frac{1}{2}$\quad (81 theories)}

At the maximum value $a/N^2 = \tfrac{1}{2}$, the $\mathcal{B}$- and $\mathcal{A}$-point coefficients vanish on \emph{both} nodes for every morphology: $\alpha_{\mathcal{B}} = \alpha_{\mathcal{A}} = 0$. The conformal window for additional fundamental flavours is entirely $N$-independent, bounded only by the $\mathcal{O}(1)$ constant $\beta$ that depends on the S/A character of the rank-2 matter. Class~44 is the endpoint of the formula $\alpha_{\mathcal{B}} = 1 - N_{\rm bif}/4$ extrapolated to $N_{\rm bif} = 4$, generalised to all morphologies.

\textit{Morphology $(0,\,0,\,6)$.}\;Six bifundamentals, no rank-2 tensors on either node. Both nodes are already at the $\mathcal{O}$-point conformal threshold ($\alpha_{\mathcal{O}} = 0$). At both $\mathcal{A}$ and $\mathcal{B}$, $\alpha = 0$ as well: no fundamental deformation is possible in the large-$N$ limit.

\textit{Morphology $(1,\,1,\,4)$.}\;One rank-2 equivalent unit per node and four bifundamentals. At $\mathcal{O}$, node 1 admits up to a constant number of fundamentals (independent of $N$) while node 2 is at the threshold. Both $\mathcal{B}$- and $\mathcal{A}$-point bounds have $\alpha = 0$: the conformal window size is $\mathcal{O}(1)$ and $N$-independent.

\textit{Morphology $(\tfrac{3}{2},\,\tfrac{3}{2},\,3)$.}\;Three rank-2 equivalent units per node with three bifundamentals. Both nodes admit at most $\mathcal{O}(1)$ fundamentals at $\mathcal{O}$, and $\alpha_{\mathcal{B}} = \alpha_{\mathcal{A}} = 0$: $N$-independence is confirmed.

\textit{Morphology $(2,\,2,\,2)$.}\;Two rank-2 equivalent units per node and two bifundamentals. At $\mathcal{O}$ some configurations of the rank-2 matter place node 2 at the conformal threshold while node 1 can admit a constant number of fundamentals. In all cases $\alpha_{\mathcal{B}} = \alpha_{\mathcal{A}} = 0$: the conformal window is $N$-independent.

\textit{Morphology $(\tfrac{5}{2},\,\tfrac{5}{2},\,1)$.}\;Five rank-2 equivalent units per node and one bifundamental. Both nodes are near the $\mathcal{O}$-point threshold and admit at most $\mathcal{O}(1)$ fundamentals. The $\alpha = 0$ property at both $\mathcal{A}$ and $\mathcal{B}$ confirms that this morphology, like all others in class~44, lies beyond the reach of large-$N$ fundamental deformations.

\begin{longtable}{l|cc|cc}
\caption{Coefficients $\alpha$ of $N$ in the upper bounds $N_f < \alpha\,N + \beta$ at each boundary point, for all morphologies of class~44. The constant $\beta$ (not shown) depends on the S/A character of the rank-2 matter. Negative $\alpha_{\mathcal{B}}$ or $\alpha_{\mathcal{A}}$ means the boundary condition is automatically satisfied at large $N$; `conf.' means the node is already at the conformal threshold at $\mathcal{O}$.}\label{tab:nf_cl44}\\
\toprule
$(N_{r2}^{(1)},\,N_{r2}^{(2)},\,N_{\rm bif})$ & $\alpha_{\mathcal{O}}^{(1)}$ & $\alpha_{\mathcal{B}}^{(1)}$ & $\alpha_{\mathcal{O}}^{(2)}$ & $\alpha_{\mathcal{A}}^{(2)}$ \\
\midrule
\endhead
$(0,\,0,\,6)$ & $\mathrm{conf.}$ & $\mathrm{conf.}$ & $\mathrm{conf.}$ & $\mathrm{conf.}$ \\
$(1,\,1,\,4)$ & $\mathrm{conf.}$ & $\mathrm{conf.}$ & $\mathrm{conf.}$ & $\mathrm{conf.}$ \\
$(\tfrac{3}{2},\,\tfrac{3}{2},\,3)$ & $\mathrm{conf.}$ & $\mathrm{conf.}$ & $\mathrm{conf.}$ & $\mathrm{conf.}$ \\
$(2,\,2,\,2)$ & $\mathrm{conf.}$ & $\mathrm{conf.}$ & $\mathrm{conf.}$ & $\mathrm{conf.}$ \\
$(\tfrac{5}{2},\,\tfrac{5}{2},\,1)$ & $\mathrm{conf.}$ & $\mathrm{conf.}$ & $\mathrm{conf.}$ & $\mathrm{conf.}$ \\
\bottomrule
\end{longtable}

\subsection*{Summary}

The results above reveal three qualitatively distinct behaviours of the conformal window under fundamental deformations, summarised in Table~\ref{tab:nf_summary}.

\begin{enumerate}

\item \textbf{Rational compression (classes 1 and 9).} Both classes have rational R-charges at the fixed point, leading to rational $\alpha_{\mathcal{B}}$ and $\alpha_{\mathcal{A}}$. For class~1, the uniform compression factor is exactly $\tfrac{1}{2}$: $\alpha_{\mathcal{B}} = 1$ against $\alpha_{\mathcal{O}} = 2$. For class~9, the compression follows the exact formula $\alpha_{\mathcal{B}} = 1 - N_{\rm bif}/4$ across all symmetric morphologies, reflecting the fixed-point bifundamental R-charge $R_{\rm bif} = \tfrac{1}{2}$ of the $a/N^2 = 27/64$ solution. The constant offset $\beta$ is identical at $\mathcal{O}$ and $\mathcal{B}$, so the compression acts only on the $N$-scaling of the conformal window.

\item \textbf{$N$-independent conformal window (class 44).} At the maximum $a/N^2 = \tfrac{1}{2}$, $\alpha_{\mathcal{B}} = \alpha_{\mathcal{A}} = 0$ for every morphology. The conformal window is bounded solely by the $\mathcal{O}(1)$ constant $\beta$, which can be at most 4 (for the $(2,2,2)$ morphology with both nodes carrying adjoint-type rank-2 tensors). No fundamental deformation with $N_f \sim N$ is possible in this class, regardless of the matter configuration. Class~44 is the endpoint $N_{\rm bif} = 4$ of the class~9 formula, extended to all morphologies.

\item \textbf{Universal tightening (classes 2, 4, 16, 21, 23).} For all classes with irrational R-charges, the bifundamental R-charges at the fixed point satisfy $R_{\rm bif} < \tfrac{2}{3}$, giving a negative anomalous dimension $\gamma_{\rm bif} = 3R_{\rm bif} - 2 < 0$. The resulting shift $\Delta\alpha = N_{\rm bif}|\gamma_{\rm bif}|/2$ tightens \emph{both} boundary conditions relative to $\mathcal{O}$: nodes with $\alpha_{\mathcal{O}} > 0$ see their coefficient reduced, while nodes at the UV threshold ($\alpha_{\mathcal{O}} = 0$) acquire $\alpha_{\mathcal{B/A}} = -\Delta\alpha < 0$, automatically satisfying their boundary condition for large $N$. The binding (positive) boundary coefficient increases monotonically with $a/N^2$: $\approx 0.219$ (class~2, $a/N^2 \approx 0.152$), $\approx 0.427$ (class~16, $a/N^2 \approx 0.443$), $\approx 0.675$ (class~21, $a/N^2 \approx 0.453$), $\approx 0.853$ (class~23, $a/N^2 \approx 0.458$), approaching unity as $a/N^2 \to \tfrac{1}{2}$. Class~4 ($a/N^2 \approx 0.168$) has the largest rank-2 content ($N_{r2}^{(1)} = \tfrac{5}{2}$) among irrational classes, giving the lighter node a wide tightened window $\alpha_{\mathcal{A}}^{(2)} \approx 1.243$. In all cases the conformal window is controlled by the $\mathcal{B}$ or $\mathcal{A}$ condition, not by the asymptotic-freedom condition at $\mathcal{O}$.

\end{enumerate}

Across all three regimes, the constant $\beta$ in the $\mathcal{B}$/$\mathcal{A}$ bounds equals the constant in the $\mathcal{O}$ bound. The IR dynamics thus affect only the large-$N$ scaling of the conformal window, leaving the $\mathcal{O}(1)$ structure --- which encodes the S/A character of the rank-2 matter --- unchanged.

\begin{table}[h]
\centering\small
\renewcommand{\arraystretch}{1.3}
\caption{Coefficients of $N$ in the three upper bounds on the number of additional fundamental flavours, for representative morphologies of each $\mathrm{SU}(N)\times\mathrm{SU}(N)$ non-Veneziano class. For asymmetric classes with a UV-saturated node ($\alpha_{\mathcal{O}}=0$), that node is listed as node~1; for class~2 (no UV-saturated node), node~1 is the lighter node. Irrational values are given to four decimal places.}
\label{tab:nf_summary}
\begin{tabular}{c|c|ccc|ccc}
\toprule
Class & $a/N^2$ & $\alpha_{\mathcal{O}}^{(1)}$ & $\alpha_{\mathcal{B}}^{(1)}$ & $\alpha_{\mathcal{A}}^{(1)}$ & $\alpha_{\mathcal{O}}^{(2)}$ & $\alpha_{\mathcal{B}}^{(2)}$ & $\alpha_{\mathcal{A}}^{(2)}$ \\
\midrule
$1$ & $0$ & $2$ & $1$ & $1$ & $2$ & $1$ & $1$ \\
$2$ & $-\frac{297}{32}+\frac{135\sqrt{5}}{32}$ & $2$ & $1.2188$ & $—$ & $1$ & $—$ & $0.2188$ \\
$4$ & $-\frac{229}{32}+\frac{65\sqrt{13}}{32}$ & $0$ & $-0.7569$ & $—$ & $2$ & $—$ & $1.2431$ \\
$9$ & $\frac{27}{64}\ (N_{\rm bif}=1)$ & $1$ & $\frac{3}{4}$ & $\frac{3}{4}$ & $1$ & $\frac{3}{4}$ & $\frac{3}{4}$ \\
$9$ & $(N_{\rm bif}=2)$ & $1$ & $\frac{1}{2}$ & $\frac{1}{2}$ & $1$ & $\frac{1}{2}$ & $\frac{1}{2}$ \\
$9$ & $(N_{\rm bif}=3)$ & $1$ & $\frac{1}{4}$ & $\frac{1}{4}$ & $1$ & $\frac{1}{4}$ & $\frac{1}{4}$ \\
$9$ & $(N_{\rm bif}=4)$ & $1$ & $0$ & $0$ & $1$ & $0$ & $0$ \\
$16$ & $\frac{3}{32}+\frac{5\sqrt{5}}{32}$ & $0$ & $-0.5730$ & $—$ & $1$ & $—$ & $0.4270$ \\
$21$ & $-\frac{7}{2}+\frac{5\sqrt{10}}{4}$ & $0$ & $-0.3246$ & $—$ & $1$ & $—$ & $0.6754$ \\
$23$ & $\frac{2295\sqrt{17}}{86528}+\frac{30159}{86528}$ & $0$ & $-0.1468$ & $—$ & $1$ & $—$ & $0.8532$ \\
$44$ & $\frac{1}{2}$ & $—$ & $0$ & $0$ & $—$ & $0$ & $0$ \\
\bottomrule
\end{tabular}
\end{table}
